\section{Related Work}
\label{sec:related}

\textbf{Cluster-level schedulers.} Shared-state and two-level schedulers
such as those described by Chen et al.~\cite{chen2015omega} and Vasquez and
Holm~\cite{vasquez2016mesoslike} decouple resource offering from placement
policy, letting frameworks bid for resources independently. Tempo operates
at the placement-policy layer and could sit underneath either architecture;
our contribution is orthogonal to the offer-vs-assign split those systems
address.

\textbf{Latency-aware placement.} Raman and Okafor~\cite{raman2020tail}
show that shared caches and network contention, not just CPU scheduling,
dominate tail latency in multi-tenant fleets, and propose isolating
latency-sensitive tenants onto dedicated node pools. This sidesteps the
packing problem rather than solving it; Tempo instead tries to make dense
co-location safe by predicting when it will hurt.

\textbf{Sampling-based load balancing.} Batch-sampling
schedulers~\cite{holm2018batch} and multi-resource packers such as
Ferro and Tan's~\cite{ferro2019tetris} extension of bin-packing heuristics
to vector resource demands both use point-in-time load samples. We compare
directly against a power-of-two-choices baseline in
Section~\ref{sec:results} and find that Tempo's predictive signal
outperforms it once utilization crosses roughly 60\%.

\textbf{Predictive autoscaling.} Tan and Ferro~\cite{tan2022predictive}
apply lightweight latency models to horizontal autoscaling decisions at
the fleet level. Tempo applies a structurally similar predictor, but at
the much finer grain of individual placement decisions rather than
fleet-wide capacity changes, and folds migration cost into the same model.

\textbf{Speculative execution and stragglers.} Raman and
Holm~\cite{raman2023straggler} address tail latency caused by individual
slow tasks via speculative re-execution. That mechanism is complementary
to Tempo: stragglers within an already-placed job are a separate problem
from deciding where to place work in the first place. Serverless-specific
placement work by Okafor et al.~\cite{okafor2020serverless} and
cluster-manager retrospectives by Okafor and Raman~\cite{okafor2017borgish}
and Ferro and Chen~\cite{ferro2021elastic} inform our choice of which
signals are cheap enough to compute inside a hot scheduling path.
